%% main.tex
%% ---------------------------------------------------------------------
%% Documento principal do template nao-oficial de monografia/dissertacao/
%% tese do IDP. Baseado no arquivo idpthesis.cls.
%%
%% Para compilar (linha de comando):
%%   pdflatex main.tex
%%   biber main
%%   pdflatex main.tex
%%   pdflatex main.tex
%%
%% Troque a opcao da classe abaixo conforme o seu trabalho:
%%   monografia | dissertacao | tese
%% ---------------------------------------------------------------------
\documentclass[monografia]{idpthesis}

%% ================== METADADOS (edite estes campos) ====================
\instituicao{Instituto Brasileiro de Ensino, Desenvolvimento e Pesquisa -- IDP}
\curso{Direito}                                  % nome do curso (graduacao)
% \programapos{Direito Constitucional}           % use p/ dissertacao/tese
% \areaconcentracao{Direito P\'ublico}           % use p/ dissertacao/tese

\tituloTrabalho{T\'itulo do Trabalho em Letras Mai\'usculas}
\subtituloTrabalho{Subt\'itulo, se houver}
\autorTrabalho{Nome Completo do Autor}

\orientador{Prof(a). Esp., Ms. (ou Dr.) Nome Completo do Orientador}
% \coorientador{Nome Completo do Coorientador}

\localdefesa{Bras\'ilia-DF}
\anotrabalho{2026}
\datadefesa{5 de agosto de 2026}

\membrosbanca{%
  \examinador{Prof. Nome Completo}{Orientador}
  \examinador{Prof. Nome Completo}{Examinador}
  \examinador{Prof. Nome Completo}{Examinador}
}

\addbibresource{referencias.bib}

\begin{document}

%% ============================================================
%% ELEMENTOS PRE-TEXTUAIS (NBR 14724, secao 4.1)
%% ============================================================
\capa
\folioderosto
\fichacatalografica
\folhadeaprovacao

\begin{dedicatoria}
  Dedico este trabalho \`a minha fam\'ilia, por todo apoio,
  confian\c{c}a e incentivo ao longo desta jornada.
\end{dedicatoria}

\begin{agradecimentos}
  Agrade\c{c}o a todos que contribu\'iram, direta ou indiretamente,
  para a realiza\c{c}\~ao deste trabalho: professores, orientador(a),
  fam\'ilia e amigos.
\end{agradecimentos}

\begin{resumo}{Palavra-chave 1. Palavra-chave 2. Palavra-chave 3.}
  Texto do resumo em l\'ingua vern\'acula, entre 150 e 500 palavras,
  contendo objetivo geral, principais conceitos e autores da
  fundamenta\c{c}\~ao te\'orica, metodologia, resultados alcan\c{c}ados
  e conclus\~oes. Consulte a NBR 6028 em caso de d\'uvidas.
\end{resumo}

\begin{resumoingles}{Keyword 1. Keyword 2. Keyword 3.}
  Abstract text in English, stating the purpose of the work, key
  theoretical concepts, methodology, main results and conclusions.
\end{resumoingles}

\listadeilustracoes
\listadetabelas
% \listofquadros
% \listofgraficos

\begin{listadeabreviaturas}
  \sigla{DPP}{Decanato de Pesquisa e P\'os-Gradua\c{c}\~ao}
  \sigla{MEC}{Minist\'erio da Educa\c{c}\~ao}
\end{listadeabreviaturas}

\sumario

%% ============================================================
%% ELEMENTOS TEXTUAIS (NBR 14724, secao 4.2)
%% A partir daqui a numeracao de p\'agina passa a ser exibida.
%% ============================================================
\idpiniciotextual

\chapter{Introdu\c{c}\~ao}
\label{cap:introducao}

Na primeira p\'agina da parte textual, apresente o seu trabalho. A
introdu\c{c}\~ao deve ocupar uma ou, no m\'aximo, duas p\'aginas. Nela,
dê um panorama da tem\'atica, exponha a quest\~ao-problema que conduzir\'a
a pesquisa proposta, a hip\'otese e a justificativa.

Os t\'itulos de cada se\c{c}\~ao (1 INTRODU\c{C}\~AO, 2 REFERENCIAL
TE\'ORICO etc.) devem ser grafados em letras mai\'usculas e negritadas,
sem ponto ap\'os a numera\c{c}\~ao (a classe \texttt{idpthesis} j\'a faz
isso automaticamente para \verb|\chapter|).

\section{Objetivo geral}
Descreva aqui o objetivo geral do trabalho.

\section{Objetivos espec\'ificos}
\begin{itemize}
  \item Objetivo espec\'ifico 1;
  \item Objetivo espec\'ifico 2.
\end{itemize}

\section{Problema}
Apresente a quest\~ao-problema que conduz a pesquisa.

\section{Hip\'otese}
Apresente a hip\'otese de pesquisa.

\section{Justificativa}
Justifique a relev\^ancia do trabalho.

Caso haja necessidade de notas explicativas, use notas de rodap\'e\footnote{Texto
de exemplo para nota de rodap\'e. A nota de rodap\'e apresenta informa\c{c}\~oes
ou coment\'arios adicionais sobre o conte\'udo.}, que a classe já formata em
fonte tamanho 10, conforme a norma.

\chapter{Referencial Te\'orico}
\label{cap:referencial}

Texto introdut\'orio de um par\'agrafo, conforme exigido pela ABNT. Esta
parte do trabalho demanda o dom\'inio das regras de cita\c{c}\~ao (diretas
e indiretas). Reveja as regras estabelecidas pela ABNT (NBR 10520) para
cita\c{c}\~oes diretas e indiretas e indica\c{c}\~ao de fontes de gr\'aficos,
figuras ou outro recurso ilustrativo utilizado.

\section{T\'itulo representativo do primeiro conceito}
\label{sec:conceito1}

Discuss\~ao espec\'ifica sobre a articula\c{c}\~ao das informa\c{c}\~oes com
foco na delimita\c{c}\~ao proposta.

\subsection{T\'itulo representativo de um subt\'opico}

Exemplo de cita\c{c}\~ao direta de at\'e 3 linhas, que permanece no corpo do
texto, entre aspas: de acordo com \textcite{sa1971}, isso acontece ``por meio
da mesma arte de conversa\c{c}\~ao que abrange t\~ao extensa e significativa
parte da nossa exist\^encia cotidiana''.

Exemplo de cita\c{c}\~ao direta com mais de 3 linhas (ambiente
\texttt{citacao}, recuo de 4\,cm, fonte tamanho 10, espa\c{c}amento simples,
sem aspas), conforme \textcite{nicholls1993}:

\begin{citacao}
A teleconfer\^encia permite ao indiv\'iduo participar de um encontro
nacional ou regional sem a necessidade de deixar seu local de origem.
Tipos comuns de teleconfer\^encia incluem o uso da televis\~ao, telefone e
computador. Atrav\'es de \'audio-confer\^encia, utilizando a companhia
local de telefone, um sinal de \'audio pode ser emitido em um sal\~ao de
qualquer dimens\~ao. \parencite[p.~181]{nicholls1993}
\end{citacao}

Exemplo de cita\c{c}\~ao indireta: para \textcite{authier1982}, a ironia
seria assim uma forma impl\'icita de heterogeneidade.

Exemplo de inser\c{c}\~ao de figura com fonte obrigat\'oria:

\begin{figure}[htbp]
  \centering
  % \includegraphics[width=0.7\textwidth]{figuras/exemplo.png}
  \fbox{\parbox{0.7\textwidth}{\centering\vspace{2cm}Espa\c{c}o reservado para a figura\vspace{2cm}}}
  \caption{T\'itulo da figura}
  \label{fig:exemplo}
  \fonte{Do autor (2026).}
\end{figure}

Para mais informa\c{c}\~oes sobre cita\c{c}\~ao, consulte a norma NBR 10520:2002.

\chapter{Procedimentos Metodol\'ogicos}
\label{cap:metodologia}

Apresente de forma clara e objetiva o modo pelo qual o projeto foi
implementado e, ap\'os um breve texto introdut\'orio, explicite a
metodologia em etapas ou a\c{c}\~oes. Poder\'a, ainda, fazer refer\^encias
\`as t\'ecnicas ou detalhes espec\'ificos utilizados.

\chapter{An\'alise e Discuss\~ao dos Dados}
\label{cap:resultados}

Apresente os resultados alcan\c{c}ados/esperados e fa\c{c}a uma an\'alise
comparativa entre os resultados esperados e os alcan\c{c}ados com a
execu\c{c}\~ao do projeto de pesquisa ou interven\c{c}\~ao.

\chapter{Considera\c{c}\~oes Finais}
\label{cap:conclusao}

Fa\c{c}a uma reflex\~ao sobre a quest\~ao-problema. Verifique se a
hip\'otese e os objetivos geral e espec\'ificos foram alcan\c{c}ados.
Apresente suas sugest\~oes e cuidados, caso outras pessoas queiram
implementar as a\c{c}\~oes propostas em suas atividades profissionais.


%% ============================================================
%% ELEMENTOS POS-TEXTUAIS (NBR 14724, secao 4.3)
%% ============================================================
\referencias

\apendices
\apendice{T\'itulo do Ap\^endice A}
\label{ape:a}

Conte\'udo do ap\^endice (documento elaborado pelo pr\'oprio autor para
complementar sua argumenta\c{c}\~ao).


\anexos
\anexo{T\'itulo do Anexo A}
\label{ane:a}

Conte\'udo do anexo (documento n\~ao elaborado pelo autor, usado para
fundamenta\c{c}\~ao ou comprova\c{c}\~ao).


\end{document}
